\section{Three Deployment Case Studies}\label{sec:discussion}
Applying our repository schema to the literature resolves apparent contradictions across three fundamental deployment decisions (Table~\ref{tab:synthesis-decisions}): (1) when upstream traffic shaping survives downstream workflow hops; (2) why read-only query decoys and live GUI humanization produce opposite task outcomes; and (3) how adaptive retraining ($D_2$) and mechanism-specific cost boundaries govern operating-point selection.

\subsection{Case 1: bind the final target-bearing emitter}
\paragraph{The downstream re-emission boundary}
When an upstream LLM API applies record padding or token batching, does that protection cover the user-facing connection after an application frontend? NetEcho demonstrates that a content-aware frontend that consumes decoded upstream tokens and re-serializes responses into new chunks (Trace~A token counts per group; Trace~B character lengths) restores topic and conversation distinguishability \cite{netecho2025}. SplitDecoy establishes the same boundary on a backward return path: even when forward split-model activations are obfuscated with decoys, the returned backward gradient exposes nonzero row support on every tested frame across nine seeds \cite[Sections 5.1--5.3]{politis2026splitdecoy}.

\noindent\textbf{Proposition 1 (target-oblivious postprocessing bound).}
Let $S$ be the hidden target, $Z$ the complete protected upstream representation and $Y$ the downstream observation. If $S\to Z\to Y$ forms a Markov chain (any downstream framing or randomization is conditionally independent of $S$ given $Z$), then by the data processing inequality \cite[Theorem~2.8.1]{cover2006elements},
\[
 I(S;Y)\leq I(S;Z).
\]
A content-aware frontend that decrypts, strips padding and re-chunks plaintext ($S\to (Z,\text{plaintext})\to Y$), or a backward pass that computes gradients from unprotected loss masks, need not satisfy $I(S;Y\mid Z)=0$. The sufficient-condition theorem then ceases to certify the bound; leakage must be checked separately. Section~\ref{sec:controlled-placement} demonstrates the specified re-emission mechanism on TLS~1.2 sockets.

\subsection{Case 2: resolve the WRA versus AHB cover asymmetry}
\paragraph{Two axes of asymmetry: threat polarity and action risk}
A comparison that groups WRA's decoy requests \cite{jeong2026wra} and AHB's humanized actions \cite{zhu2026turing} together as ``cover traffic'' can create an apparent paradox: WRA reports a clean privacy improvement (one concurrent decoy reduces functional-intent reconstruction $E_{\rm func}$ from 0.770 to 0.610), whereas AHB reports severe task collapse (Trip Planning success falls from 0.75 Raw to 0.15 under combined \texttt{Online+Fake} humanization).

Binding their result records shows that WRA and AHB sit at opposite poles along \emph{two structural axes}:
\begin{enumerate}
\item \textbf{Threat polarity inversion (Type~I vs.\ Type~III):} WRA is a Type~I privacy defense protecting a user's research intent from a network/domain observer. AHB is a Type~III bot-evasion mechanism designed to make an automated mobile GUI agent mimic human touch patterns so it bypasses platform anti-automation controls.
\item \textbf{Action-interface risk magnitude (discarded read-only queries vs.\ live GUI state mutations):} WRA's decoys are \emph{read-only search and retrieval prompts} whose returned pages are discarded locally before synthesis; they never invoke state-mutating tools on the user's account. By contrast, AHB's \texttt{Fake} operator injects \emph{unfiltered live GUI touch events} directly onto stateful mobile application screens, triggering unintended navigation and form submissions. The non-Fake \texttt{History+Long} configuration changes timing/trajectory history. These source-reported condition rates are not treated as a matched effect: condition-specific sample denominators were not independently recovered in this revision.
\end{enumerate}

\paragraph{Dispatch-time authorization versus terminal-state scoring}
Separating cover execution routes also exposes a fundamental \emph{analytical design argument} regarding agent benchmark evaluators: terminal state-snapshot scorers (such as $\tau$-bench's native reward, which hashes the final database state against the reference state) are \textbf{path-oblivious}. If an unconfirmed live cover action happens to perform the exact state transition that the user's task later intended (or preempts a real action so the real action fails with a status error while leaving the terminal state identical to baseline), a terminal-snapshot scorer awards a perfect score of $1.0$ although the endpoint does not establish authorization of the earlier call. Systems admitting cover to live dispatch must therefore enforce authorization and provenance checks at dispatch time and inspect tool state-machine transitions (such as one-way status flags or append-only audit ledgers, demonstrated in Section~\ref{sec:tau-cover-bridge}).

\subsection{Case 3: bind adaptive retraining ($D_2$) and cost objects}
\paragraph{The $D_0\to D_2$ rebound and component-cost boundaries}
Selecting a defense operating point requires pairing an adaptive attacker evaluation (comparing a clean-trained classifier $D_0$ against a defended-trained classifier $D_2$ on the \emph{exact same defended test set}) with the cost of the actual deployed mechanism. In UI Traces \cite{uitraces2026}, retraining on timestamp-delayed archive traces ($D_2$) recovers substantial model-attribution accuracy (confirmed by our supplementary archive reanalysis on 1,135 test records: $+0.1503$ nats log-loss reduction over $D_0$, 98.333\% Bonferroni-adjusted CI (three contrasts) $[0.1205,0.1815]$; Supplement S4). Section~\ref{sec:controlled-placement} gives a matched TLS adaptation example. The separate retail study shows why prior-based baselines and rule attacks must accompany a learned observer; an improvement over $D_0$ alone does not establish useful inference.

On the cost side, record binding prevents cross-mechanism price substitution: Early Bird's $6$ms/$4\%$ GPTCache semantic de-identification overhead does not price $K$-prefix KV-cache isolation \cite{earlybird2025}; KVGov's $1.3\,\mu$s HMAC key-derivation microcost omits workload-dependent lost cache reuse \cite{addagada2026kvgov}; and SpliceDefense's $<5$ms/$<1\%$ QCP padding overhead prices chunk padding rather than combined CTBF+QCP deployment \cite{spliceleak2026}. Conversely, CachePrune jointly binds marked-token attack suppression, QA/summarization utility and serving throughput on selective KV sharing \cite{wu2026cacheprune}, and OTRO bounds conditional access-count leakage ($3.94\to 0.04$ bits given byte length) alongside separately scoped unpadded TTFT ($4.59\%$) and epoch-padding ($2.8$ms) costs \cite{lee2026otro}.
\begin{table*}[t]\centering\footnotesize
\caption{Source-local findings, cross-study synthesis and the resulting measurement obligation. These are reuse judgments, not accusations that the source authors claimed end-to-end protection.}\label{tab:synthesis-decisions}
\begin{tabularx}{\textwidth}{@{}p{.14\textwidth}YYY@{}}\toprule
Check & Source-local evidence and locator & Synthesis added by record separation & Changed judgment / controlled evidence\\\midrule
1. Receiver & NetEcho IV-C--E/Fig.5: frontend re-emission; SplitDecoy 5.1--5.3: backward gradient support. & Protected upstream and forward views do not cover newly exposed downstream or backward inputs. & Bind the final observer or prove the postprocessing condition. Two-link TLS placement reproduces this mechanism; costs use identical texts.\\
2. Execution & WRA 8.1/Table8: read-only query decoys; AHB 5.2.3/Table1/E.3: live GUI humanization. & Objectives, live action semantics and protected state scopes differ, so their utility outcomes are not contradictory. & Evaluate task objects, full account and shared database separately. Retail same-item writes change status; known-position filtering defeats cover in the specified dispatch view.\\
3. Adaptation and cost & UI Traces 6.1: delayed-trace refit; Dangol 9.7: capture versus live protocol; Early Bird's KV-prefix versus semantic-cache mechanisms. & A paired attack contrast requires the same defended tests; costs from another mode or mechanism cannot price the claimed protection. & Keep source-local comparisons and mark missing bridges unresolved. TLS adapted accuracy remains 70.8\% under 1024B buckets; the retail class-prior baseline exceeds the legacy learned accuracy.\\\bottomrule
\end{tabularx}\end{table*}

